\chapter{Special and Limit Solutions}

\section*{1. Generalized fixed points}

The Ricci flow has only a small class of genuine fixed points. On a compact
manifold, the unnormalized flow fixes exactly Ricci-flat metrics, while the
normalized flow fixes exactly Einstein metrics. The larger class of
self-similar solutions, also called Ricci solitons, behaves as generalized
fixed points under diffeomorphism and scaling.

\begin{remark}
\label{ch02-rem-self-similar-solutions}
\source{chow-knopf-ricci-flow:page-0034,chow-knopf-ricci-flow:page-0035}
\dcref{Remark 2.1}
\group{generalized-fixed-points}

By rescaling in the soliton equation, one may assume that the parameter
\(\lambda\) is one of \(-1,0,1\). These cases correspond respectively to
shrinking, steady, and expanding solitons.
\end{remark}

\begin{remark}
\label{ch02-rem-gradient-solitons}
\source{chow-knopf-ricci-flow:page-0035}
\dcref{Remark 2.2}
\group{generalized-fixed-points}

If the soliton vector field is a gradient field \(X=-\nabla f\), then the
soliton equation becomes \(\nabla\nabla f=\Rc+\lambda g\), and the metric is
called a gradient Ricci soliton.
\end{remark}

\begin{example}
\label{ch02-ex-gaussian-soliton}
\source{chow-knopf-ricci-flow:page-0036}
\dcref{Example 2.3}
\group{generalized-fixed-points}

Euclidean space with its standard metric is both a steady Ricci soliton and,
with the potential function \(f(x)=\frac12|x|^2\), an expanding gradient Ricci
soliton. In the latter form it is called the Gaussian soliton.
\end{example}

\begin{lemma}
\label{ch02-lem-self-similar-solutions-and-solitons}
\source{chow-knopf-ricci-flow:page-0036}
\dcref{Lemma 2.4}
\group{generalized-fixed-points}

A solution of the Ricci flow is self-similar if and only if its initial metric
is a Ricci soliton.
\end{lemma}

\begin{proof}
\source{chow-knopf-ricci-flow:page-0036}
\sketch If \(g(t)=\sigma(t)\psi_t^*g_0\) is self-similar, then differentiating at
\(t=0\) gives
\[
-2\Rc(g_0)=\sigma'(0)g_0+\mathcal L_{Y(0)}g_0,
\]
where \(Y(t)\) generates \(\psi_t\). This is exactly the soliton equation
\(-2\Rc=\mathcal L_X g_0+2\lambda g_0\) with \(X=Y(0)\) and
\(\lambda=\frac12\sigma'(0)\).

Conversely, if \(g_0\) satisfies the soliton equation for some complete vector
field \(X\), define \(\sigma(t)=1+2\lambda t\) and let \(\psi_t\) be generated by
the time-dependent vector field \(Y_t=\sigma(t)^{-1}X\). Then
\[
g(t)=\sigma(t)\psi_t^*g_0
\]
satisfies \(\partial_t g=-2\Rc[g]\) by the standard pullback and Lie-derivative
calculation.
\end{proof}

\section*{2. Eternal solutions}

An eternal solution is defined for all \(t\in(-\infty,\infty)\).

\subsection*{2.1. The cigar soliton}

\begin{example}
\label{ch02-ex-cigar-soliton}
\source{chow-knopf-ricci-flow:page-0037,chow-knopf-ricci-flow:page-0038}
\dcref{Example 2.4}
\group{eternal-solutions}

Hamilton's cigar soliton is the complete surface \((\R^2,g_\Sigma)\) with
\[
g_\Sigma=\frac{dx\otimes dx+dy\otimes dy}{1+x^2+y^2}
=\frac{dr^2+r^2\,d\theta^2}{1+r^2}
=ds^2+\tanh^2(s)\,d\theta^2,
\]
scalar curvature \(R_\Sigma=\frac{4}{1+r^2}=\frac{4}{\cosh^2 s}\), and area form
\(\frac{1}{1+r^2}\,dx\wedge dy\). It is asymptotic to a cylinder of radius \(1\)
with exponential decay.
\end{example}

\begin{lemma}
\label{ch02-lem-cigars-unique}
\source{chow-knopf-ricci-flow:page-0039,chow-knopf-ricci-flow:page-0040,chow-knopf-ricci-flow:page-0041}
\dcref{Lemma 2.7}
\group{eternal-solutions}

Up to homothety, the cigar soliton is the unique rotationally symmetric
gradient Ricci soliton of positive curvature on \(\R^2\).
\end{lemma}

\begin{proof}
\source{chow-knopf-ricci-flow:page-0039,chow-knopf-ricci-flow:page-0040,chow-knopf-ricci-flow:page-0041}
\sketch Write a rotationally symmetric metric as \(g=ds^2+\varphi(s)^2\,d\theta^2\).
Using the orthonormal frame \(e_1=\partial_s\), \(e_2=\varphi^{-1}\partial_\theta\),
the connection satisfies \(\omega^2_1=\varphi'(s)\,d\theta\), so the curvature is
\(K=-\varphi''/\varphi\). For a radial potential \(f=f(s)\), the steady gradient
soliton equation \(Kg=\nabla\nabla f\) reduces to
\[
f''=-\frac{\varphi''}{\varphi},\qquad
\frac{\varphi' f'}{\varphi}=-\frac{\varphi''}{\varphi}.
\]
Hence \(f'=\alpha\varphi\) for a constant \(\alpha\neq 0\), and convexity forces
\(\alpha>0\). Substituting back gives \(\varphi''+2a^2\varphi\varphi'=0\), so
\(\varphi'+a^2\varphi^2=b\). Smoothness at the origin requires
\(\varphi(0)=0\) and \(\varphi'(0)=1\), forcing \(b=1\). Thus
\(\varphi(s)=a^{-1}\tanh(as)\), and after the rescaling \(\sigma=as\) one gets
\(g=a^{-2}g_\Sigma\).
\end{proof}

\begin{remark}
\label{ch02-rem-cigars-potential}
\source{chow-knopf-ricci-flow:page-0041}
\dcref{Remark 2.8}
\group{eternal-solutions}

For the cigar soliton, the potential may be taken as
\(f(s)=a^{-1}\log(\cosh(as))\), and the corresponding soliton vector field is
\(X=-\nabla f=-\tanh(as)\,\partial_s\).
\end{remark}

\subsection*{2.2. The Rosenau solution}

\begin{lemma}
\label{ch02-lem-cylinder-to-sphere}
\source{chow-knopf-ricci-flow:page-0042,chow-knopf-ricci-flow:page-0043}
\dcref{Lemma 2.10}
\group{eternal-solutions}

A rotationally symmetric metric
\[
g=\varphi(z)^2\,dz^2+\psi(z)^2\,g_{\mathrm{can}}
\]
on \((-\omega,\omega)\times S^n\) extends smoothly to \(S^{n+1}\) if and only if
\(\int_{-\omega}^{\omega}\varphi<\infty\), \(\psi(\pm\omega)=0\),
\(\psi'(\pm\omega)/\varphi(\pm\omega)=\mp 1\), and all even derivatives of
\(\psi\) with respect to arc length vanish at the endpoints.
\end{lemma}

\begin{proof}
\source{chow-knopf-ricci-flow:page-0042,chow-knopf-ricci-flow:page-0043}
\sketch The endpoint conditions are exactly those ensuring that, after changing to arc
length \(s\) and then to Euclidean coordinates \((x,y)\) near a pole, the metric
has a smooth even expansion in the radial variable. The vanishing of the
specified derivatives is equivalent to smooth extendibility across the pole.
\end{proof}

\begin{lemma}
\label{ch02-lem-rosenau-sphere}
\source{chow-knopf-ricci-flow:page-0046,chow-knopf-ricci-flow:page-0047}
\dcref{Lemma 2.12}
\group{eternal-solutions}

The Rosenau metric
\[
u(x,t)=\frac{2\beta\sinh(-\alpha\lambda t)}{\cosh(\alpha x)+\cosh(\alpha\lambda t)}
\]
defines an ancient solution on \(S^2\) if and only if \(\alpha=1\) and
\(\beta=\frac{1}{2\lambda}\).
\end{lemma}

\begin{proof}
\source{chow-knopf-ricci-flow:page-0046,chow-knopf-ricci-flow:page-0047}
\sketch The conformal-factor evolution equation reduces the Ricci flow condition to the
identity \(2\beta\lambda=\alpha\). The smooth extension across the poles is
checked using Lemma~2.10 with \(\varphi=\psi=\sqrt{u}\): the distance to the
poles is finite, the conformal factor vanishes at the ends, and the slope
condition forces \(\alpha=1\). The higher endpoint conditions follow from the
explicit decay of derivatives. Substituting then gives \(\beta=\frac1{2\lambda}\).
\end{proof}

\begin{remark}
\label{ch02-rem-rosenau-limit-cigar}
\source{chow-knopf-ricci-flow:page-0047}
\dcref{Exercise 2.13}
\group{eternal-solutions}

The Rosenau solution is Type II and, when viewed infinitely far back in time
near either pole, limits to the cigar soliton.
\end{remark}

\section*{4. Immortal solutions}

An immortal solution exists for all sufficiently large times \(t>\alpha\). The
example below expands from a cone and becomes a smooth complete metric of
positive curvature for \(t>0\).

\begin{lemma}
\label{ch02-lem-expanding-cone-soliton}
\source{chow-knopf-ricci-flow:page-0048,chow-knopf-ricci-flow:page-0049,chow-knopf-ricci-flow:page-0050}
\dcref{Lemma 2.15}
\group{immortal-solutions}

There exists an expanding self-similar Ricci-flow solution on \(\R^2\) of the
form \(g(t)=t(f(r)^2\,dr^2+r^2\,d\theta^2)\) whose profile \(f\) satisfies
\[
f'=\frac r2(f^2-f^3),
\]
equivalently
\[
\frac{r^2}{4}=C-\frac1{f(r)}+\log\left|\frac{f(r)}{f(r)-1}\right|,
\]
and whose positive-curvature branch is given explicitly by the Lambert-\(W\)
formula in the text.
\end{lemma}

\begin{proof}
\source{chow-knopf-ricci-flow:page-0048,chow-knopf-ricci-flow:page-0049,chow-knopf-ricci-flow:page-0050}
\sketch For the ansatz \(g(t)=t(f(r)^2dr^2+r^2d\theta^2)\), direct computation of the
Christoffel symbols and Ricci tensor shows that the modified flow equation
\(\partial_t g=-2\Rc+\mathcal L_X g\), with
\(X=\frac{r}{tf(r)}\partial_r\), is equivalent to the single ODE
\(f'=\frac r2(f^2-f^3)\). This ODE integrates by partial fractions to the
implicit relation above. The curvature formula
\(K=\frac1{2t}\frac{1-f}{f}\) shows that positive curvature requires
\(0<f<1\); on that branch, rewriting in terms of \(F=f^{-1}-1\) and inverting
\(Fe^F\) with the Lambert-\(W\) function gives the stated explicit solution.
\end{proof}

\section*{5. The neckpinch}

\begin{theorem}
\label{ch02-thm-neckpinch-rigorous}
\source{chow-knopf-ricci-flow:page-0053}
\dcref{Theorem 2.15}
\group{neckpinch}

For \(n\ge 2\), an open set of \(\SO(n+1)\)-invariant metrics on \(S^{n+1}\)
develops a Type-I neckpinch at a finite time \(T\), and any parabolic blow-up
at the developing singularity converges to the shrinking cylinder soliton
\(ds^2+2(n-1)(T-t)g_{\mathrm{can}}\).
\end{theorem}

\begin{theorem}
\label{ch02-thm-neckpinch-asymptotics}
\source{chow-knopf-ricci-flow:page-0053}
\dcref{Theorem 2.16}
\group{neckpinch}

Let \(\bar s(t)\) be the location of the smallest neck. Then near a neckpinch
there are inner- and intermediate-layer estimates comparing \(\psi(x,t)\) to
\(\sqrt{2(n-1)(T-t)}\), with quadratic control in the inner layer and a logarithmic
control in the intermediate layer.
\end{theorem}

\begin{remark}
\label{ch02-rem-single-point-pinching}
\source{chow-knopf-ricci-flow:page-0054}
\dcref{Remark 2.17}
\group{neckpinch}

For a reflection-symmetric dumbbell with one neck and two bumps, the
neckpinch can occur only on the totally geodesic hypersurface at the neck,
provided the diameter stays bounded as the singular time is approached.
\end{remark}

\begin{remark}
\label{ch02-rem-neckpinch-pde}
\source{chow-knopf-ricci-flow:page-0054}
\dcref{Remark 2.18}
\group{neckpinch}

The \(s\)-derivative and time derivative do not commute; rather,
\([\partial_t,\partial_s]=-n(\psi_{ss}/\psi)\partial_s\).
\end{remark}

\begin{remark}
\label{ch02-rem-neckpinch-curvatures}
\source{chow-knopf-ricci-flow:page-0054}
\dcref{Remark 2.19}
\group{neckpinch}

For the warped-product metric \(ds^2+\psi^2g_{\mathrm{can}}\), the sectional
curvatures are \(K_0=-\psi_{ss}/\psi\) and \(K_1=(1-\psi_s^2)/\psi^2\), and
the Ricci tensor is \(\Rc=(nK_0)ds^2+(K_0+(n-1)K_1)g_{\mathrm{can}}\).
\end{remark}

\begin{proposition}
\label{ch02-prop-neckpinch-bounds}
\source{chow-knopf-ricci-flow:page-0055}
\dcref{Proposition 2.20}
\group{neckpinch}

For a rotationally symmetric solution with \(|\psi_s|\le 1\) and \(R>0\)
initially, the slope bound persists, the curvature is controlled by
\(\psi^{-2}\), the minimum of \(K_1\) is nondecreasing, the scalar curvature
stays positive, \(\psi_t<0\), \(\psi^2\) is uniformly Lipschitz in time, and
\(\psi(x,T)\) exists pointwise if the solution exists up to time \(T\).
\end{proposition}

\begin{lemma}
\label{ch02-lem-neckpinch-slope}
\source{chow-knopf-ricci-flow:page-0056}
\dcref{Lemma 2.21}
\group{neckpinch}

If \(|\psi_s|\le 1\) holds initially and the boundary conditions at the poles
hold, then the slope \(v=\psi_s\) satisfies \(1\le \sup|v(\cdot,t)|\le \sup|v(\cdot,0)|\).
\end{lemma}

\begin{lemma}
\label{ch02-lem-neckpinch-a-evolution}
\source{chow-knopf-ricci-flow:page-0056}
\dcref{Lemma 2.22}
\group{neckpinch}

The scale-invariant quantity \(a=\psi^2(K_1-K_0)\) evolves by
\[
a_t=a_{ss}+(n-4)\frac{\psi_s}{\psi}a_s-4(n-1)\frac{\psi_s^2}{\psi^2}a.
\]
\end{lemma}

\begin{corollary}
\label{ch02-cor-neckpinch-a-bound}
\source{chow-knopf-ricci-flow:page-0057}
\dcref{Corollary 2.23}
\group{neckpinch}

The quantity \(a\) remains uniformly bounded by its initial supremum:
\(\sup|a(\cdot,t)|\le \sup|a(\cdot,0)|\).
\end{corollary}

\begin{corollary}
\label{ch02-cor-neckpinch-curvature-radius}
\source{chow-knopf-ricci-flow:page-0057}
\dcref{Corollary 2.24}
\group{neckpinch}

There is a constant \(C\) depending only on \(n\) and the initial metric such
that \(|\Rm|\le C/\psi^2\).
\end{corollary}

\begin{lemma}
\label{ch02-lem-neckpinch-L-evolution}
\source{chow-knopf-ricci-flow:page-0057}
\dcref{Lemma 2.25}
\group{neckpinch}

The curvature quantity \(L=K_1\) evolves by
\[
L_t=\Delta L+2\frac{\psi_s}{\psi}L_s+2\bigl[K^2+(n-1)L^2\bigr].
\]
\end{lemma}

\begin{corollary}
\label{ch02-cor-neckpinch-Lmin}
\source{chow-knopf-ricci-flow:page-0058}
\dcref{Corollary 2.26}
\group{neckpinch}

The minimum of \(L\) is nondecreasing; if \(L_{\min}(0)\neq 0\), then
\[
L_{\min}(t)\ge \frac{1}{L_{\min}^{-1}(0)-2(n-1)t}.
\]
\end{corollary}

\begin{lemma}
\label{ch02-lem-neckpinch-shrinking-radius}
\source{chow-knopf-ricci-flow:page-0058}
\dcref{Lemma 2.27}
\group{neckpinch}

If \(|\psi_s|\le 1\) initially and the scalar curvature is initially positive,
then positivity of scalar curvature persists and \(\psi_t<0\) for as long as
the solution exists.
\end{lemma}

\begin{lemma}
\label{ch02-lem-neckpinch-lipschitz}
\source{chow-knopf-ricci-flow:page-0058,chow-knopf-ricci-flow:page-0059}
\dcref{Lemma 2.28}
\group{neckpinch}

Under the same hypotheses, \(|(\psi^2)_t|\le 2(\alpha+n)\), where
\(\alpha=\sup|a(\cdot,0)|\).
\end{lemma}

\begin{corollary}
\label{ch02-cor-neckpinch-limit-profile}
\source{chow-knopf-ricci-flow:page-0059}
\dcref{Corollary 2.29}
\group{neckpinch}

If the solution exists on \(0\le t<T\), then \(\psi(x,t)\) has a pointwise
limit as \(t\nearrow T\) for each \(x\in[-1,1]\).
\end{corollary}

\begin{proposition}
\label{ch02-prop-neckpinch-profile}
\source{chow-knopf-ricci-flow:page-0059,chow-knopf-ricci-flow:page-0060}
\dcref{Proposition 2.30}
\group{neckpinch}

For a rotationally symmetric solution with \(|\psi_s|\le 1\) and \(R>0\)
initially, the number of necks cannot increase, the smallest neck shrinks at a
controlled Type-I rate, the solution is concave on the polar cap, and if the
right-most bump persists with limit height \(D>0\), then no singularity occurs
on the polar cap.
\end{proposition}

\begin{lemma}
\label{ch02-lem-neckpinch-zeroes}
\source{chow-knopf-ricci-flow:page-0060}
\dcref{Lemma 2.31}
\group{neckpinch}

The slope \(v=\psi_s\) has finitely many zeroes, the number of zeroes is
nonincreasing in time, and the number drops if a degenerate critical point
appears.
\end{lemma}

\begin{lemma}
\label{ch02-lem-neckpinch-neck-rate}
\source{chow-knopf-ricci-flow:page-0060,chow-knopf-ricci-flow:page-0061}
\dcref{Lemma 2.32}
\group{neckpinch}

If \(|\psi_s|\le 1\) and \(R>0\) initially, then the smallest neck radius
satisfies
\[
(n-1)(T-t)\le r_{\min}(t)^2\le 2(n-1)(T-t),
\]
so the solution must become singular by time \(r_{\min}(0)^2/(n-1)\) unless it
loses all necks first.
\end{lemma}

\begin{lemma}
\label{ch02-lem-tensor-maximum-boundary}
\source{chow-knopf-ricci-flow:page-0061}
\dcref{Lemma 2.33}
\group{neckpinch}

For a compact manifold with boundary, if a symmetric tensor evolves by
\(\partial_t S=\Delta S+U*S\) and is positive initially and on the boundary, and
if the reaction term preserves the null cone of \(S\), then \(S\) remains
nonnegative.
\end{lemma}

\begin{lemma}
\label{ch02-lem-neckpinch-concavity-polar-cap}
\source{chow-knopf-ricci-flow:page-0061,chow-knopf-ricci-flow:page-0062}
\dcref{Lemma 2.34}
\group{neckpinch}

If \(\psi_{ss}(x,0)\le 0\) on the polar cap, then \(\psi_{ss}(x,t)\le 0\) on
the polar cap for all later times.
\end{lemma}

\begin{lemma}
\label{ch02-lem-neckpinch-no-polar-singularity}
\source{chow-knopf-ricci-flow:page-0062,chow-knopf-ricci-flow:page-0063,chow-knopf-ricci-flow:page-0064,chow-knopf-ricci-flow:page-0065}
\dcref{Lemma 2.35}
\group{neckpinch}

If the right-most bump persists to a positive limiting height \(D>0\), then no
singularity occurs on the polar cap beyond the neck.
\end{lemma}

\begin{proposition}
\label{ch02-prop-neckpinch-typei-cylinder}
\source{chow-knopf-ricci-flow:page-0066,chow-knopf-ricci-flow:page-0067,chow-knopf-ricci-flow:page-0068,chow-knopf-ricci-flow:page-0069,chow-knopf-ricci-flow:page-0070}
\dcref{Proposition 2.36}
\group{neckpinch}

For a maximal rotationally symmetric solution with \(|\psi_s|\le 1\), \(R\ge 0\),
and at least one neck, the smallest neck becomes singular in finite time, the
singularity is Type I, a scale-invariant pinching estimate holds for \(K/L\),
and the neck profile is asymptotic to the shrinking cylinder with
\(r_{\min}(t)\sim \sqrt{2(n-1)(T-t)}\).
\end{proposition}

\begin{lemma}
\label{ch02-lem-neckpinch-typei-bound}
\source{chow-knopf-ricci-flow:page-0066}
\dcref{Lemma 2.37}
\group{neckpinch}

Under the hypotheses of Proposition~2.36, the full curvature satisfies
\(|\Rm|\le C/(T-t)\).
\end{lemma}

\begin{lemma}
\label{ch02-lem-neckpinch-F-evolution}
\source{chow-knopf-ricci-flow:page-0066,chow-knopf-ricci-flow:page-0067}
\dcref{Lemma 2.38}
\group{neckpinch}

The quantity \(F=\frac{K}{L}\log L\) satisfies a parabolic evolution equation
with lower-order terms \(P\) and \(Q\) as in the source.
\end{lemma}

\begin{lemma}
\label{ch02-lem-neckpinch-Fhat-bound}
\source{chow-knopf-ricci-flow:page-0067,chow-knopf-ricci-flow:page-0068}
\dcref{Lemma 2.39}
\group{neckpinch}

The scale-invariant quantity
\[
\hat F=\frac{K}{L}\bigl[\log L+2-\log L_{\min}(0)\bigr]
\]
satisfies \(\sup_{\mathcal W(t)}\hat F\le \max\{n-1,\sup_{\mathcal W(0)}\hat F\}\).
\end{lemma}

\begin{lemma}
\label{ch02-lem-neckpinch-neck-profile-estimates}
\source{chow-knopf-ricci-flow:page-0068,chow-knopf-ricci-flow:page-0069,chow-knopf-ricci-flow:page-0070}
\dcref{Lemma 2.40}
\group{neckpinch}

Near the smallest neck, \(\psi/r_{\min}\) satisfies inner- and
intermediate-layer estimates of the form displayed in the source, with the
same logarithmic scaling.
\end{lemma}

\begin{proposition}
\label{ch02-prop-neckpinch-happen}
\source{chow-knopf-ricci-flow:page-0071,chow-knopf-ricci-flow:page-0072,chow-knopf-ricci-flow:page-0073}
\dcref{Proposition 2.42}
\group{neckpinch}

There exist smooth rotationally symmetric initial metrics on \(S^{n+1}\) with
\(|\Psi_s|\le 1\), positive scalar curvature, and a sufficiently small neck so
that the Ricci flow develops a neckpinch in finite time.
\end{proposition}

\begin{lemma}
\label{ch02-lem-neckpinch-bump-model}
\source{chow-knopf-ricci-flow:page-0071}
\dcref{Lemma 2.43}
\group{neckpinch}

For \(W(s)=\sqrt{A+Bs^2}\) with the stated range of \(B\), the metric
\(ds^2+W(s)^2g_{\mathrm{can}}\) has positive scalar curvature and pinching
quantity \(a\) satisfying \(1-B<a<1+B\).
\end{lemma}

\begin{lemma}
\label{ch02-lem-neckpinch-initial-family}
\source{chow-knopf-ricci-flow:page-0072,chow-knopf-ricci-flow:page-0073}
\dcref{Lemma 2.44}
\group{neckpinch}

There is a smoothed family of initial metrics \(g_{A,B}\) with positive scalar
curvature, slope bound \(|\psi_s|\le 1\), a neck of radius \(\sqrt A\), a bump
of height at least \(A\sqrt B\), and uniformly bounded pinching quantity.
\end{lemma}

\begin{proposition}
\label{ch02-prop-single-point-pinching}
\source{chow-knopf-ricci-flow:page-0074}
\dcref{Proposition 2.45}
\group{neckpinch}

If the diameter remains bounded as the singular time is approached, then the
solution stays positive away from the equatorial hypersurface, so the
singularity occurs only at the neck.
\end{proposition}

\begin{lemma}
\label{ch02-lem-single-point-subsolution}
\source{chow-knopf-ricci-flow:page-0074,chow-knopf-ricci-flow:page-0075}
\dcref{Lemma 2.46}
\group{neckpinch}

If the limiting growth parameter \(\rho(T)\) is finite, then the comparison
function \(v=\psi_s\) dominates the linear subsolution \(\underline v\) on the
interior region up to time \(T\).
\end{lemma}

\begin{remark}
\label{ch02-rem-cigar-negative-curvature-exercise}
\dcref{Exercise 2.9}
\group{eternal-solutions}
\source{chow-knopf-ricci-flow:page-0041}
The negative-curvature analog of the cigar is obtained by solving the same
rotationally symmetric soliton ODE with \(K<0\).
\end{remark}

\begin{remark}
\label{ch02-rem-degenerate-neckpinch-picture}
\dcref{Problem 2.47}
\group{degenerate-neckpinch}
\source{chow-knopf-ricci-flow:page-0075,chow-knopf-ricci-flow:page-0076,chow-knopf-ricci-flow:page-0077,chow-knopf-ricci-flow:page-0078}
The expected transition between neckpinch and spherical extinction suggests
that a parameter value exists at which the Ricci flow forms a Type IIa
singularity, i.e. a slowly forming singularity.
\end{remark}

\begin{remark}
\label{ch02-rem-rosenau-singularity-model}
\dcref{Remark 2.14}
\group{eternal-solutions}
\source{chow-knopf-ricci-flow:page-0047}
The Rosenau solution is relevant as a possible dimension-reduction limit of a
three-manifold singularity, although Perelman's no-local-collapsing theorem
rules out that scenario for finite-time singularities.
\end{remark}

\begin{remark}
\label{ch02-rem-slowly-forming-meaning}
\dcref{Notes and commentary}
\group{degenerate-neckpinch}
\source{chow-knopf-ricci-flow:page-0079}
No example of a slowly forming singularity on a compact manifold was known at
the time of the source; a rigorous example on \(\R^2\) had recently been
constructed.
\end{remark}

\begin{remark}
\label{ch02-rem-typei-typeiia}
\dcref{Remark 2.47}
\group{degenerate-neckpinch}
\source{chow-knopf-ricci-flow:page-0077,chow-knopf-ricci-flow:page-0078}
A finite-time singularity is called Type I, or rapidly forming, when
\((T-t)\sup|\Rm|\) stays bounded; it is called Type IIa, or slowly forming, when
\((T-t)\sup|\Rm|\) is unbounded.
\end{remark}
